% !TeX root = ../papers/nc-p3.tex
\section{同時文字交換が残る二直線と二つの印を持つ正則窓}
\label{sec:flip-character-locus}
$S_z\left(\begin{smallmatrix}a&b\\c&d\end{smallmatrix}\right)
=\left(\begin{smallmatrix}d&c/z\\zb&a\end{smallmatrix}\right)$とする。
\begin{theorem}\label{thm:simultaneous-character-flip}
規格化した可逆チャートで、四位置全てで$i\mapsto i+1$を誘導する単項自己同型が
存在する必要十分条件は、ある同じ$\sigma=\pm1$について
$S_z A_0=\sigma A_0$, $S_z A_1=\sigma A_1$となることである。
第一行列の条件$d=\sigma a$, $c=\sigma zb$だけで第二行列の条件も従う。
自己同型の平方は各位置で、射影的には低ブロック上の1と高ブロック上の$\sigma$である。
従って二ブロックと三候補を復元する前節までの条件の下では、
$\sigma=+1$の直線を除いて文字分解は一意である。この直線では実際に二つの異なる
印を持つ反対階数16の標準正則窓が存在する。
\end{theorem}
\begin{proof}
$u^2=z$とし、同時文字交換の後、四位置の成分ごとに順に
\[
 (u,u,1,z),\quad(1/z,1,1/u,1/u),\quad
 (u,u,1,z),\quad(1,1/z,1/u,1/u)
\]
を掛ける。薄い成分は元に戻る。厚い上左ブロックは二行列とも$S_z$で変わり、
二つの交差ブロックは不変である。右下はSchur補で一意に定まるため、
正規形は$(S_zA_0,S_zA_1,z,\lambda)$となる。
対角変換でこの二行列をさらに同じ定数だけ掛けられる。
逆に規格化を保つ対角変換の残りを確認する。薄い二色では低群から高群への
$2\times2$が全て非零である。二色で右・左薄ベクトルに許される逆倍率を
$t_2,t_3$とすると、低群の二行と高群の二列でたすき掛けに比較して
$t_3/t_2=\epsilon=\pm1$である。高群の二行・低群の二列でも、薄い元係数を
左右から消せば同じ比となる。従って四位置とも各組内の比は$\epsilon$である。
位置0の低・高の代表倍率を$A,B$、位置2を$C,D$とすると、薄い式で位置1,3の
倍率も決まり、厚い交差ブロックを保つ条件は$AD=BC$となる。
二つの上左ブロックには同じ$(A/B)^2$が掛かる。
特に対角自己同型なら$A/B=\pm1$で、位置ごとのスカラーを除けば同時文字符号だけである。
$S_z^2=1$なので、二行列を同時倍率で戻せるのは定理の固有条件に限られる。
第一固有条件を補完式に代入すると$g_0=zf_0$, $h_0=e_0$で、
負の固有値では両定数が零になる。また$g_1=\sigma z f_1$, $h_1=\sigma e_1$であり、
任意の根に対して$g=\sigma zf$, $h=\sigma e$が従う。

$u^2=\sigma z$と取り直した同じ四倍率と文字交換は、固有条件上の窓を固定する。
各二文字でこの単項変換の平方は二つの倍率の積であり、高・低ブロックの
その比は全位置で$\sigma$である。対角文字符号を掛けても積は両ブロックに
同一の符号だけ掛かり、この比は変わらない。
別の文字分解があれば、その下位ビット符号は三候補の非自明な別元である。
各二文字の射影直線で元の符号と可換で固定点を共有しないため、元の二文字を交換する。
四位置への伝播で全てのブロックに同じ結論が出る。これは射影位数2である必要があり、
上の必要条件は正の直線に絞られる。

具体的に$(z,\lambda)=(3,2)$で
\[
 A_0=\begin{pmatrix}0&1\\3&0\end{pmatrix},\qquad
 A_1=\begin{pmatrix}2&8/13\\24/13&2\end{pmatrix}
\]
とする。$x=-192/169,y=1$で二つの符号は$-1,1$、従って既約な対根族上にある。
三切断の階数は$6,16,6$で、四次中心次元は$(3,2,0,0)$である。
四中心切断による商は素数1009で全始点・次数13まで$(1+t+t^2+t^3)^4$の
次元を持つ。定理\ref{thm:kummer-ramification}と同じ上下界と自由性の証明により
この有理窓は標準Koszul左右AS正則である。

同時文字交換の各倍率に下位ビット符号を掛け、低・高の各二文字でその二固有直線を
新しい文字に取る。各位置で固有値の組は低・高で等しく、総テンソルは不変なので、
新しい文字も総次数零となる。これらの直線はいずれも元の座標直線と異なる。
二つの切断で厚い各ブロックの四つの二行二列主・交差ブロックが全て可逆で、
薄い二色の全成分が非零であることを、$\sqrt3,\sqrt{-1}$を含む良い還元1009で確認する。
両切断の全階数6とこれらの下界から各色の階数は$2,2,1,1$となる。
$z$に対応する薄い二行列の零点対は元と同じ交差比を与えるので、新表示の
$z$も$3$または$1/3$である。\path{computations/verify_flip_character_locus.py}
は倍率、二直線の係数、四中心切断および新固有基底の非零条件を検算する。
\end{proof}
